\bookchapter{Глава 51. Разработка интерфейса — знакомство с Jetpack Compose}

\chapterintro{Интерфейс — всё, что видит пользователь. До появления Compose интерфейсы Android писали на XML — муторно, как заполнять пропуски в упражнениях. В 2021 году Google выпустила Jetpack Compose: интерфейс «рисуют» прямо кодом на Kotlin, и результат виден сразу — поменялась строка кода, предпросмотр мгновенно обновился. К 2026 году Compose стал безоговорочным мейнстримом для новых проектов. В этой главе мы научимся «рисовать».}

\booksection{Идея Compose: интерфейс — это функция}

Ядро идеи Compose — одна фраза: интерфейс = функция. Всё, что может появиться на экране (текст, кнопка, картинка), — «композируемая функция» (Composable), помеченная @Composable. Вложите их друг в друга, словно кубики, — вот вам и интерфейс.

\begin{lstlisting}[language=Kotlin]
@Composable
fun Greeting(name: String) {
    Text(text = "Привет, $name!")     // «кирпичик текста»
}

// Вложенная компоновка: большие блоки содержат маленькие
@Composable
fun HomeScreen() {
    Column {                          // вертикально
        Greeting(name = Иван)
        Greeting(name = Мария)
        Button(onClick = { /* что делать по нажатию */ }) {
            Text("Нажми")
        }
    }
}
\end{lstlisting}

Обратите внимание на «шляпу» @Composable — она говорит компилятору: «это кубик интерфейса». Кубики можно вкладывать без ограничений — в этом весь секрет Compose.

\booksection{Базовые компоненты: текст, кнопка, поле ввода, изображение}

Компоненты Compose — те же «кубики». Вот четыре самых ходовых:

\begin{center}
\begin{tabularx}{\linewidth}{|X|X|X|}
\hline
\textbf{Компонент} & \textbf{Что делает} & \textbf{Обязательные параметры} \tabularnewline
\hline
Text & Показывает текст & text (текст для показа) \tabularnewline
Button & Кнопка & onClick (колбэк нажатия) \tabularnewline
TextField & Поле ввода & value, onValueChange (колбэк изменения ввода) \tabularnewline
Image & Показывает изображение & painter (графический ресурс) \tabularnewline
\hline
\end{tabularx}
\end{center}

\begin{center}\small\color{gray}Четыре базовых компонента Compose\end{center}

\begin{lstlisting}[language=Kotlin]
@Composable
fun BasicWidgets() {
    Column {
        Text(text = "Это какой-то текст")
        Button(onClick = { println("Нажато") }) {
            Text("Кнопка")
        }
        var input by remember { mutableStateOf("") }
        TextField(
            value = input,
            onValueChange = { input = it },
            label = { Text("Введите что-нибудь") }
        )
    }
}
\end{lstlisting}

В двух строках про поле ввода мелькнуло remember \{ mutableStateOf(“”) \} — это «магия состояния» Compose; подробно разберём в пятом разделе главы. Пока просто запомните её в лицо.

\booksection{Трио раскладок: Column, Row, Box}

Компоненты есть — как их расставить? Запомните три «коробки»:

\begin{itemize}
  \item Column: ставит элементы столбиком (как очередь за обедом);
  \item Row: выстраивает элементы в ряд (как сидящие рядом);
  \item Box: кладёт элементы стопкой (как колода карт; удобно для бейджей и наложений).
\end{itemize}

\begin{lstlisting}[language=Kotlin]
Column {
    Text("Первая строка")
    Row {
        Text("Слева")
        Text("Справа")
    }
    Box {
        Text("Нижний слой")
        Text("Бейдж поверх")
    }
}
\end{lstlisting}

Хотите, чтобы раскладка выглядела лучше? Добавьте компонентам Modifier (модификатор) — он словно «накладывает грим» на кубики:

\begin{lstlisting}[language=Kotlin]
Text(
    text = "Красивый текст",
    modifier = Modifier
        .padding(16.dp)          // отступы со всех сторон
        .fillMaxWidth()          // занять всю ширину
        .background(Color.Blue)  // цвет фона
        .clickable { /* нажатие */ }// кликабельно
)
// Цепочки вызовов: одна за другой, порядок важен (сначала padding, потом фон — фон не включает отступы)
\end{lstlisting}

\begin{tcbwarn}{Внимание: порядок в Modifier имеет значение}
padding(16.dp).background(Blue) и background(Blue).padding(16.dp) дают разный результат: первый — «сначала отступ, потом заливка» (фон не захватывает отступы), второй — «сначала заливка, потом отступ» (фон включает отступы). Новички чаще всего рисуют здесь «загадочные белые поля». Сомневаетесь — поменяйте порядок, и предпросмотр тут же даст ответ.

\end{tcbwarn}

\booksection{Состояние: оживляем интерфейс}

Статичный интерфейс — скучное зрелище. Чтобы после нажатия кнопки число росло на 1, а в поле ввода можно было печатать, нужно состояние (State). Правило Compose: интерфейс всегда верно отражает состояние; состояние изменилось — интерфейс сам перерисовался.

\begin{lstlisting}[language=Kotlin]
@Composable
fun Counter() {
    // remember + mutableStateOf: запоминаем меняющееся значение
    var count by remember { mutableStateOf(0) }

    Column {
        Text(text = "Вы нажали $count раз")
        Button(onClick = { count++ }) {   // меняем состояние
            Text("Нажми +1")                // интерфейс обновится сам
        }
    }
}
\end{lstlisting}

Три строки кода — и живой счётчик. Вот чем Compose приятнее старых подходов: не нужно вручную «обновлять интерфейс» — достаточно поменять данные. Запомните две ключевые функции: remember (не даёт значению обнулиться при рекомпозиции интерфейса) и mutableStateOf (сообщает Compose: «это значение изменилось, надо перерисовать»).

Пользователь нажимает

Состояние меняется (count++)

Compose замечает изменение

Интерфейс перерисовывается сам (рекомпозиция)

\begin{center}\small\color{gray}Рисунок 3. Однонаправленный поток данных «состояние интерфейс» в Compose\end{center}

\booksection{Тема и тёмный режим: настроил один раз — работает везде}

MyFirstAppTheme \{ ... \} из шаблонного проекта — это и есть тема. Она построена на дизайн-системе Material 3 и помогает единообразно управлять цветами, шрифтами и формами. Поддержка тёмного режима — всего три шага:

\begin{enumerate}
  \item определите в теме две палитры — светлую и тёмную (шаблон уже генерирует ColorScheme автоматически);
  \item когда система переходит в тёмный режим, isSystemInDarkTheme() определяет это автоматически;
  \item берите цвета из темы (например, MaterialTheme.colorScheme.primary), а не вписывайте значения напрямую.
\end{enumerate}

\begin{tcbinfo}{Совет: динамические цвета Material 3}
Начиная с Android 12 «Material You» позволяет приложению брать цвета из ваших обоев. Шаблон Compose поддерживает это из коробки (dynamicColor). Ваше приложение будет «переодеваться» на телефонах разных пользователей — вот она, романтика современного Android.

\end{tcbinfo}

\booksection{Предпросмотр, анимации и живое редактирование: «читы» Compose}

Три «чита» Compose для удобства разработки: попробуешь — назад дороги нет:

\begin{itemize}
  \item Предпросмотр @Preview: добавьте Composable аннотацию @Preview — и Android Studio мгновенно отрисует интерфейс справа: правите код и сразу видите результат, без установки на телефон каждый раз;
  \item Живое редактирование (Live Edit): Android Studio 2026 года умеет «горячо» применять часть правок на эмуляторе или реальном устройстве — мелкие изменения вроде цвета текста появляются на устройстве за секунды;
  \item API анимаций: animateFloatAsState и AnimatedVisibility позволяют парой строк добиться плавных переходов элементов — до ощущения «дорогого» продукта часто не хватает именно этих нескольких десятков строк.
\end{itemize}

\begin{lstlisting}[language=Kotlin]
@Composable
fun AnimatedBox() {
    var big by remember { mutableStateOf(false) }
    // при изменении значения плавно переходит к целевому
    val size by animateDpAsState(if (big) 120.dp else 60.dp)

    Box(Modifier
        .size(size)
        .background(MaterialTheme.colorScheme.primary)
        .clickable { big = !big }
    )
}
\end{lstlisting}

\begin{tcbinfo}{Совет: сперва функциональность, потом анимация}
Типичная болезнь новичков — «интерфейс ещё не сделан, а человек уже мучается с анимациями». Сначала добейтесь, чтобы всё работало, а затем разом добавьте анимацию — это и есть принцип «простота превыше всего»: анимация — украшение, а функциональность — спасение.

\end{tcbinfo}

\begin{tcbpractice}{Практика: сделайте «кнопку настроения»}
Превратите шаблонный проект в переключатель настроения: одна кнопка; каждое нажатие двигает настроение по кругу радость грусть злость, а текст и цвет фона меняются вслед за ним. Подсказка: выбирайте текст и цвет через when по состоянию. Готовым приложением поделитесь с друзьями — и соберёте первый урожай «Вау, ты уже умеешь делать приложения!».

\end{tcbpractice}

\begin{tcbsummary}{Коротко о главном}
\begin{itemize}
  \item Идея Compose: интерфейс = функции @Composable, вложенные как кубики
  \item Четыре главных компонента: Text / Button / TextField / Image
  \item Трио раскладок: Column — столбиком, Row — в ряд, Box — стопкой
  \item Modifier «накладывает грим» цепочкой: padding, fillMaxWidth, background, clickable — порядок важен
  \item Состояние — душа интерфейса: remember + mutableStateOf; поменял данные — интерфейс перерисовался сам
  \item Тема: Material 3 + динамические цвета; тёмный режим настраивается один раз и работает везде
\end{itemize}

\end{tcbsummary}
